\section{Conclusion}
\label{chp:conclusion}

% This dissertation argues that the mismatch between modern hardware capabilities and
% legacy memory management designs is the root cause of critical \emph{performance}
% bottlenecks---and that targeted, hardware- and workload-aware changes to those systems
% eliminate the bottlenecks without compromising security or extensibility.

This dissertation argues that 
designing memory management systems with deep understanding of hardware capabilities and workload characteristic 
can eliminate performance bottlenecks without compromising security or extensibility.

We identified three manifestations of the bottlenecks and addressed each with a targeted system.
First, address translation has become a major performance bottleneck on terabyte-scale
servers, yet commodity OS kernels remain hardwired to radix-tree page tables, blocking the
faster translation architectures that would relieve it.
\EMTbrand{} gives Linux an extensible, architecture-neutral interface for memory
translation, enabling OS support for radically different page table structures without
modifying the kernel for each new design; using it, we port ECPT and FPT---two recent
hardware schemes for fast translation---to Linux and reveal OS-level challenges invisible
to hardware-only evaluation.
Second, strict I/O memory protection in virtualized environments incurs 83--95\% throughput
degradation from per-DMA VM exits, forcing production deployments to disable it entirely.
\vFreebrand{} redesigns the Linux IOMMU invalidation protocol around one-for-many
invalidations and deferred free page, amortizing VM exit costs while preserving strict
security guarantees, and achieves performance within 1--8\% of an unprotected system---a
74--147$\times$ improvement over the state of the art.
Third, GPU memory oversubscription for AI workloads is unaddressed by the CUDA driver,
leaving each framework to hand-write its own offloading logic or suffer a more than
5$\times$ penalty from default UVM.
\pytorchUVMbrand{} integrates a UVM-aware caching allocator and framework-guided eviction
directly into PyTorch, enabling practical oversubscription for LLM inference without
modifying application code; under oversubscription it is 45$\times$--58$\times$ faster than
default UVM and 27$\times$--35$\times$ faster than hand-written KV-cache offloading.

Across all three, a common structure emerges.
Each begins with a mature, widely deployed abstraction---Linux's radix-tree paging, the
Linux IOMMU subsystem, and CUDA Unified Virtual Memory---that is correct and general but
was designed under assumptions that no longer hold at scale.
Closing the resulting semantic gap does not require replacing the existing stack: in each
case the fix is surgical, identifying the specific invariant that the legacy design
enforces unnecessarily and relaxing it while retaining the surrounding system.
Realizing this required cross-layer evaluation, because each bottleneck is created not by
hardware but by a software layer---the OS overhead that simulation models assume constant,
the IOMMU datapath that serializes VM exits triggered by invalidation, 
and lack of framework support for UVM.

Today's computing stacks pair increasingly heterogeneous hardware with increasingly
demanding workloads, and memory management systems must evolve to bridge them---serving
the needs of modern workloads while fully exploiting the capabilities of modern hardware.
This dissertation demonstrates that redesigning the memory management systems to be hardware- and workload-aware is a high-leverage path
to unlocking the full potential of modern hardware.